\section{The Deployed System and Its Learning Gate}
\label{sec:live}

Everything above is computed from satellites and simulation. The people being flooded know things
neither can see. The deployed system closes that loop under strict honesty constraints, and we
report the loop's status accurately, which is that its machinery works and its evidence does not
yet exist.

\subsection{Citizen reports as data}
Dashboard users drop a pin with a water level they can actually judge, being ankle, knee, waist or
chest, mapped to depth bands with explicit tolerances of $\pm$0.15 to 0.35\,m, plus an optional
note. This is why the labels in Fig.~\ref{fig:teaser} are words rather than numbers. A person can
confirm or contradict ``knee-deep''. Nobody can confirm 0.47\,m.

Reports are validated against the area's model window and rate-limited by salted address hash, with
five per hour, a per-cell cool-down, a global daily cap and a honeypot field for bots. They are held
in memory for the live map and mirrored by a background scheduler to a public dataset repository,
which doubles as durable storage across free-tier restarts. No personal data is served, and the
hashes never leave the backend.

\subsection{From reports to supervision}
A report is only a label together with the storm that caused it. The nightly job groups reports by
calendar day in local time, tags each day with the rainfall that actually fell, and drops dry days
below 5\,mm. A knee-deep report without rain is far more likely to be mischief or a burst pipe than
a storm signal. Same-cell same-day reports collapse to their median, so one street corner cannot
dominate.

\subsection{The reward gate}
Citizen reports carry no dry labels, so a model that predicts deep water everywhere would score
perfectly on wet points. That is the same failure mode as the emulator's sparse targets, now in the
field, and it is why the loop needs a gate rather than a learning rate.

A candidate emulator is fine-tuned with a band-tolerant loss on report points, meaning zero gradient
inside the reported band, anchored by a simulation replay buffer in every batch. It ships only if
all four of the following hold. There are at least 15 held-out points from at least two distinct
storm days, held out \emph{by day} because points within a day are correlated. Held-out point error
improves by at least 5\,\%. Replay error degrades by at most 10\,\%, which is the all-wet guard. And
it wins at least 70\,\% of bootstrap resamples of the held-out errors.

Every decision, accepted or held, with all four numbers, is appended to a public learning log
surfaced on the dashboard, so the loop is auditable by the people whose reports feed it.

\textbf{Status, stated plainly.} The loop is deployed and the gate is tested offline. The refresh
half, which pulls weather and recomputes outlooks, runs end to end. The learning half has never
fired, because the reports dataset is empty. Running it nightly today would produce a log of nights
with zero model updates, which is the correct behaviour and not a result. We therefore claim a
mechanism and an audit trail here, and no learning outcome. The first accepted or correctly refused
field update remains the most interesting experiment we have not been able to run.

\subsection{Live context layers}
The same discipline applies to the ambient layers. Live forecast alerts recompute the outlook
read-only and never mutate committed artifacts. Advisories are generated in English, Hindi and
Marathi from only the system's own numbers, with a deterministic template fallback so the panel
works with no language model at all. A Black Marble night lights overlay \cite{roman2018} marks
built-up cells whose latest good-quality night radiance dropped more than 50\,\% below their 90-day
median, which is a proxy for a power outage.

One trap there is worth recording for anyone building the same layer. During monsoon cloud cover a
fully quality-masked composite downloads as \emph{zero radiance}, which reads naively as a city-wide
blackout. Missing data must be treated as unknown and not as dark.
